% Options for packages loaded elsewhere
\PassOptionsToPackage{unicode}{hyperref}
\PassOptionsToPackage{hyphens}{url}
%
\documentclass[
  11pt,
]{article}
\usepackage{lmodern}
\usepackage{amssymb,amsmath}
\usepackage{ifxetex,ifluatex}
\ifnum 0\ifxetex 1\fi\ifluatex 1\fi=0 % if pdftex
  \usepackage[T1]{fontenc}
  \usepackage[utf8]{inputenc}
  \usepackage{textcomp} % provide euro and other symbols
\else % if luatex or xetex
  \usepackage{unicode-math}
  \defaultfontfeatures{Scale=MatchLowercase}
  \defaultfontfeatures[\rmfamily]{Ligatures=TeX,Scale=1}
\fi
% Use upquote if available, for straight quotes in verbatim environments
\IfFileExists{upquote.sty}{\usepackage{upquote}}{}
\IfFileExists{microtype.sty}{% use microtype if available
  \usepackage[]{microtype}
  \UseMicrotypeSet[protrusion]{basicmath} % disable protrusion for tt fonts
}{}
\makeatletter
\@ifundefined{KOMAClassName}{% if non-KOMA class
  \IfFileExists{parskip.sty}{%
    \usepackage{parskip}
  }{% else
    \setlength{\parindent}{0pt}
    \setlength{\parskip}{6pt plus 2pt minus 1pt}}
}{% if KOMA class
  \KOMAoptions{parskip=half}}
\makeatother
\usepackage{xcolor}
\IfFileExists{xurl.sty}{\usepackage{xurl}}{} % add URL line breaks if available
\IfFileExists{bookmark.sty}{\usepackage{bookmark}}{\usepackage{hyperref}}
\hypersetup{
  hidelinks,
  pdfcreator={LaTeX via pandoc}}
\urlstyle{same} % disable monospaced font for URLs
\usepackage[margin=1in]{geometry}
\setlength{\emergencystretch}{3em} % prevent overfull lines
\providecommand{\tightlist}{%
  \setlength{\itemsep}{0pt}\setlength{\parskip}{0pt}}
\setcounter{secnumdepth}{-\maxdimen} % remove section numbering
\usepackage{mathrsfs}
\usepackage{esint}
\setmathfont{Asana Math}
\setlength{\emergencystretch}{3em}

\author{}
\date{23 September 2026}

\begin{document}

\hypertarget{exact-negative-einstein-spectral-degeneration-on-non-virtually-kuxe4hler-manifolds}{%
\section{Exact negative-Einstein spectral degeneration on
non-virtually-Kähler
manifolds}\label{exact-negative-einstein-spectral-degeneration-on-non-virtually-kuxe4hler-manifolds}}

23 September 2026. This note proves a global escape alternative for the
allowed branch. It does not assert that small Weyl splitting always
vanishes. The main argument combines the previously proved signed
harmonic-Weyl estimate with classical smooth Einstein compactness and an
actual smooth-topological obstruction. Its scope is all exact outputs
satisfying the stated geometric bounds, not a neighbourhood of a
selected approximate construction.

\textbf{Strengthening for the actual aspherical family.} The subsequent
\href{line1_aspherical_no_bubbles.md}{aspherical no-bubble theorem}
removes the opposite-Weyl bound from the self-gluing application below.
In that family, vanishing splitting forces diameter escape alone.
Theorem 1 and the general degeneration discussion remain useful for
other non-virtually-Kähler types where asphericity is not available.

\hypertarget{hypotheses-and-the-new-conclusion}{%
\subsection{1. Hypotheses and the new
conclusion}\label{hypotheses-and-the-new-conclusion}}

Let \(M\) be a closed connected oriented smooth four-manifold for which
no connected cover of degree at most two admits a Kähler structure
compatible with the lifted orientation. Suppose there is \(A>0\) such
that each metric under consideration has an actual polarized monopole
class \(a_g\) with \((a_g^+)^2\ge A\). The class may depend on the
metric. The scalar monopole estimate, not an assumption about a formal
cohomology vector, is required.

Consider exact metrics \(\operatorname{Ric}_g=-g\). Order the
eigenvalues of \(W_g^+\) and set \[
 \lambda_1<\lambda_2\le\lambda_3,\qquad
 t=-\frac34\lambda_1,\quad b=\lambda_3-\lambda_2,\quad
 z=\frac b{4t}.
\] Fix \(m>0\) and \(0\le\theta<1/\sqrt3\), and impose \[
                 t\ge m,\qquad z\le\theta.                 \tag{1}
\] Define the normalized splitting energy \[
                      \mathscr S(g)=\fint_M b^2.
\]

\textbf{Theorem 1 (global bounded-geometry splitting gap).} For every
\(D,K>0\) there is \(\varepsilon=\varepsilon(A,m,\theta,D,K)>0\) such
that every metric above satisfying \[
              \operatorname{diam}(M,g)\le D,
              \qquad \|W_g^-\|_\infty\le K                 \tag{2}
\] satisfies \[
                         \mathscr S(g)\ge\varepsilon.       \tag{3}
\] The constant can be chosen uniformly over all smooth types satisfying
the same hypotheses and the same positive budget \(A\). No independent
volume, injectivity-radius, full-curvature or Sobolev hypothesis is
needed.

\textbf{Corollary 2 (escape alternative).} On a fixed such \(M\), if
exact negative-Einstein metrics satisfy (1) and \(\mathscr S(g_j)\to0\),
then \[
       \operatorname{diam}(M,g_j)+\|W^-_{g_j}\|_\infty
                         \longrightarrow\infty.            \tag{4}
\] More precisely, every fixed bounded rectangle in these two quantities
contains only metrics with a positive lower splitting energy. It is not
claimed that each summand in (4) tends to infinity separately.

This conclusion changes a global decision at the canonical shear ratio
\(T/A=1\): the earlier monopole splitting bound proportional to
\((1-T/A)^2\) vanishes there, while (3) is strictly positive under (2).
It does not exclude arbitrary outputs escaping (2).

\hypertarget{derivation-of-the-geometric-compactness-with-no-redundant-inputs}{%
\subsection{2. Derivation of the geometric compactness, with no
redundant
inputs}\label{derivation-of-the-geometric-compactness-with-no-redundant-inputs}}

Since \(s=-4\), the actual scalar monopole inequality gives \[
 16\operatorname{Vol}(g)=\int_M s^2
       \ge32\pi^2(a_g^+)^2\ge32\pi^2A,
 \qquad \operatorname{Vol}(g)\ge2\pi^2A.                    \tag{5}
\] The maximum-principle result in
\href{line1_einstein_spectral_application.md}{the Einstein spectral
application} gives \[
        t\le M_\theta=\frac1{2(1-3\theta^2)}.              \tag{6}
\] Moreover, \[
 |W^+|^2=\frac83t^2+\frac12b^2
       \le\left(\frac83+8\theta^2\right)M_\theta^2.        \tag{7}
\] Consequently the bound on \(W^-\), together with the Einstein
equation, gives a full sectional-curvature bound depending only on
\(\theta,K\).

With \(\operatorname{Ric}=-g\), Bishop comparison and the diameter bound
also give an upper volume bound depending only on \(D\). In addition,
relative Bishop comparison and (5) give a uniform lower bound for the
volume of every ball of any fixed positive radius at most \(D\). The
sectional bound and these bounds imply the standard uniform
injectivity-radius estimate. Thus every sequence satisfying (1)--(2)
has, after a subsequence and pull-back by orientation-preserving
diffeomorphisms, a smooth Einstein limit on a closed smooth manifold.
Here smoothness follows by elliptic bootstrapping of the exact Einstein
equation after Cheeger--Gromov compactness. This is classical
compactness, not a new regularity theorem.

\hypertarget{proof-of-the-splitting-gap}{%
\subsection{3. Proof of the splitting
gap}\label{proof-of-the-splitting-gap}}

If (3) failed, take a sequence with \(\mathscr S(g_j)\to0\) and use the
preceding smooth compactness. The ordered Weyl eigenvalues converge
uniformly. In the limit, \(b_\infty=0\) everywhere and
\(t_\infty\ge m\). Thus the lowest Weyl eigenline remains smooth and
separated from its two-dimensional complement.

One can already finish with the exact Derdziński theorem. Equivalently,
the proved \href{line1_signed_harmonic_weyl.md}{signed harmonic-Weyl
estimate} gives, for the normalized lowest eigenform on its orientation
cover, \[
 \fint\left(|dt_j|^2+(t_j-\tfrac12)^2+|\nabla\omega_j|^2\right)
                        \le C_{m,\theta}\mathscr S(g_j).    \tag{8}
\] The line bundles converge under the diffeomorphisms: the uniform
spectral gap makes their spectral projections smoothly convergent, and
sufficiently close real line bundles are isomorphic. Their orientation
covers are therefore identified with the limit orientation cover. This
is a cover of degree at most two, not an assumed global choice of
eigenform on \(M\).

Passing to the smooth limit in (8) gives \(t_\infty=1/2\) and a parallel
self-dual form \(|\omega_\infty|^2=2\). The corresponding orthogonal
almost complex structure is parallel, hence integrable and Kähler. The
lifted metric is negative Kähler--Einstein. Since smooth convergence
preserves the diffeomorphism type, this contradicts the topological
hypothesis for every sufficiently late manifold of the sequence. This
also proves the claimed uniformity across smooth types. Corollary 2 is
its contrapositive.

\hypertarget{actual-application-and-removal-of-the-lower-amplitude-assumption}{%
\subsection{4. Actual application and removal of the lower amplitude
assumption}\label{actual-application-and-removal-of-the-lower-amplitude-assumption}}

For the self-gluings \[
 M=M(q,g;u,f),\qquad
 (\phi,x)\longmapsto(-\phi+u(x),f(x)),\quad [u]\ne0,
 \quad q\ge1,\quad g\ge2,\quad f\in\operatorname{Diff}^+(\Sigma_g),
\] the inherited \href{round11_kahler.md}{virtual Kähler theorem}
excludes all finite Kähler covers. The
\href{round12_global.md}{canonical-class calculation} supplies an actual
canonical monopole class \(a\) with \[
 a^2=2\chi(M)>0,\qquad (a_g^+)^2\ge a^2.
\] Thus Theorem 1 applies with \(A=2\chi(M)\), to every negative
Einstein metric in its stated spectral class, independently of the
original gluing coordinates or of its position in the Einstein moduli
space.

For these objects \(\tau=0\) and \(T=A\). Hence the previously proved
topological splitting estimate involving \((1-T/A)^2\) is zero. The
present positive bound instead uses the obstruction to a degree-two
Kähler cover. It is therefore not an algebraic restatement of that
earlier estimate. No assertion is made that this compactness corollary
is new in the published literature.

The known almost-Einstein sequences on these manifolds do not contradict
the theorem: they are not exact Einstein and their long thin geometry
leaves the compact regime. The result still permits a non-Kähler exact
Einstein metric of appreciable Weyl splitting.

In fact this actual family allows removal of the lower amplitude bound.
This strengthening uses both its signature and its fundamental group,
not an assumption that the selected eigenline cannot disappear.

\textbf{Theorem 3 (shear-family gap without lower Weyl amplitude).} Fix
\(A_0>0\), \(D,K>0\), and \(0\le\theta<1/\sqrt3\). For every member of
the shear family with \(2\chi(M)\ge A_0\), every exact
\(\operatorname{Ric}_g=-g\) metric satisfying \[
 b\le4\theta t,\qquad \operatorname{diam}(g)\le D,
 \qquad\|W^-\|_\infty\le K
\] has \(\mathscr S(g)\ge\varepsilon(A_0,\theta,D,K)>0\). Here ordered
eigenvalues define \(t,b\) everywhere, and \(W^+=0\) is allowed. There
is no prescribed \(m>0\). In particular, any exact sequence in this
relative pinching class with vanishing splitting energy must satisfy the
diameter/opposite-Weyl escape conclusion (4). The later
\href{line1_aspherical_no_bubbles.md}{Theorem 3} strengthens this to
\(\operatorname{diam}(g_j)\to\infty\), with a uniform splitting gap
depending only on \(D,\theta\), and without assuming \(K\).

\textbf{Proof.} At any positive maximum of \(t\), the lowest eigenvalue
is simple because \(\lambda_2-\lambda_1=2t-b/2\ge2(1-\theta)t>0\). The
previous maximum-principle calculation therefore proves (6) even when
\(W^+\) has zeros. If \(t\equiv0\), (6) is immediate. Thus all the
compactness conclusions of Section 2 remain available, uniformly using
\(A_0\) in (5).

If a gap failed, the smooth limit would be a compact negative Einstein
metric with \(\lambda_2=\lambda_3\) everywhere. Derdziński's Proposition
5(iv) and Theorem 2 give two relevant alternatives: \(W^+\equiv0\), or a
Kähler--Einstein metric on a cover of degree at most two. The possible
non-Kähler conformally Kähler case of that classical classification has
a nonzero Killing field; a compact negative Einstein metric has none by
the integrated Bochner formula. Thus it does not provide another case
here. The KE alternative contradicts the virtual Kähler obstruction.

In the ASD alternative, the inherited signature \(\tau(M)=0\) and the
signature formula give \(\int|W^-|^2=0\). Both Weyl halves vanish, so
the limit has constant sectional curvature \(-1/3\). But the injected
vertex group \(\pi_1(F_{q,2})\times\pi_1(\Sigma_g)\) contains
\(\mathbb Z^2\). This contradicts Preissmann's theorem for closed
strictly negatively curved manifolds. Both alternatives are excluded on
every smooth type in the limiting subsequence. This proves the theorem.
In particular the argument controls the possible vanishing of the lowest
Weyl amplitude instead of hiding it in a nondegeneracy hypothesis.

Under the original lower-amplitude hypothesis, the spectral assumptions
themselves also force the endpoint volume if the splitting tends to
zero; endpoint volume is not a premise of Theorem 1. Einstein
Chern--Weil gives \[
 \frac{2\pi^2T}{V}=\frac13+\fint|W^+|^2
     =1+\frac83\fint(t^2-\tfrac14)+\frac12\mathscr S.
                                                               \tag{10}
\] By (6) and (8), the last two terms are bounded in absolute value by
\(C\sqrt{\mathscr S}+\mathscr S/2\). When \(T=A\), (5) also gives
\(V\ge2\pi^2T\). Consequently, for sufficiently small splitting, \[
 0\le\frac{V}{2\pi^2T}-1\le C_{m,\theta}\sqrt{\mathscr S}.
                                                               \tag{11}
\] Thus the hypotheses cover the full volume window before showing that
a vanishing-splitting sequence must approach its lower end. They do not
exclude candidates whose splitting is bounded away from zero.

\hypertarget{opposite-weyl-degeneration-without-the-aspherical-input}{%
\subsection{5. Opposite-Weyl degeneration without the aspherical
input}\label{opposite-weyl-degeneration-without-the-aspherical-input}}

There is a more precise necessary degeneration picture in the general
fixed-topology, bounded-diameter case. This is no longer a remaining
possibility for the aspherical self-gluings: the
\href{line1_aspherical_no_bubbles.md}{compact-core obstruction} excludes
their bubbles directly, without proving compatibility of complex
fillings. Suppose \(g_j\) are exact Einstein, satisfy (1),
\(\mathscr S(g_j)\to0\), and have bounded diameter, but no uniform
\(W^-\) bound.

The volume lower bound is still (5). The four-dimensional Einstein
Chern--Gauss--Bonnet identity gives \[
 8\pi^2\chi(M)=\int_M\left(|W^+|^2+|W^-|^2+\frac23\right),
 \qquad \operatorname{Vol}(g_j)\le12\pi^2\chi(M).           \tag{9}
\] Thus the full curvature energy is bounded and relative Bishop
comparison gives noncollapse. Classical Einstein orbifold compactness
supplies an Einstein orbifold limit with smooth convergence off finitely
many points, together with the usual ALE rescaling analysis at
concentrations.

Because (9) bounds total volume, the unnormalized splitting integral
also tends to zero. On the regular part of the limit, (8) gives \[
 t_\infty=\tfrac12,\qquad \nabla\omega_\infty=0.
\] The lowest eigenline extends as an orbifold eigenline since its
eigenvalue stays simple and nonzero. Its orientation orbifold cover
therefore carries the induced negative Kähler--Einstein orbifold metric.

At a curvature blow-up scale \(\rho_j\to0\), the rescaled metrics are
\(\rho_j^{-2}g_j\). Their Ricci constant tends to zero, while (7) gives
\[
       \|W^+_{\rho_j^{-2}g_j}\|_\infty
                     \le C_\theta\rho_j^2\longrightarrow0.
\] Consequently any smooth nonflat ALE bubble obtained in this process
is Ricci-flat and anti-self-dual in the original orientation. Wright's
classification shows that such smooth bubbles are hyperkähler or
suitable cyclic quotients of Gibbons--Hawking spaces, and are Kähler.

This is \textbf{not} a proof that the original \(M\) is Kähler. A Kähler
orbifold and Kähler ALE pieces do not by themselves provide compatible
complex structures across the actual marked necks, a global complex
smoothing, or the original smooth type. Intermediate bubble-tree objects
can themselves be orbifolds. No assertion that all bubbles are simply
connected, nor that a global complex smoothing is unobstructed, is used.
Thus, in the absence of an additional topological argument,
opposite-chirality concentration remains a genuine mechanism; it is not
silently removed by the word ``resolution''. On the actual aspherical
self-gluings, the later compact-core argument supplies precisely such an
additional argument and rules it out.

The lower amplitude hypothesis is used in this specific KE-orbifold
description. Without that hypothesis, an orbifold limit in a general
family can instead have \(W^+\equiv0\). One may not transfer
\(\tau(M)=0\) directly to its orbifold signature and repeat the
smooth-limit hyperbolicity argument: bubble and boundary correction
terms intervene. The smooth gap of Theorem 3 is unaffected.

\hypertarget{why-the-universal-small-splitting-vanishing-theorem-remains-open-here}{%
\subsection{6. Why the universal small-splitting vanishing theorem
remains open
here}\label{why-the-universal-small-splitting-vanishing-theorem-remains-open-here}}

The signed identity in the preceding work equates positive connection
energy to positive splitting energy. It does not force either to vanish.
At the negative KE curvature model the splitting reaction has
coefficient \(-6\), so there is no immediate positive-mass absorption
proof.

Near a fixed negative KE metric, Dai--Wang--Wei, Theorem 3.6, prove that
nearby Einstein metrics remain KE when all infinitesimal complex
deformations are integrable. This is a genuine sufficient geometric
condition, not an invertibility assumption. It does not give a uniform
statement over degenerating moduli or over obstructed complex
structures. Nagy's 2026 second- and third-order deformation results are
finite-order results; they cannot be substituted for all-order local
integrability.

The present theorem avoids that unresolved local issue by using actual
non-virtual-Kähler topology. It excludes smooth compact limits on those
objects regardless of the deformation theory of any hypothetical KE
limit. Diameter escape is not removed by either this note or its
aspherical strengthening. The marked ALE compatibility problem is not
solved here in general; the later aspherical result bypasses it for the
actual self-gluings by proving that an ALE core cannot embed at all.

\hypertarget{primary-sources-and-provenance}{%
\subsection{7. Primary sources and
provenance}\label{primary-sources-and-provenance}}

\begin{itemize}
\tightlist
\item
  A. Derdziński, \emph{Self-dual Kähler manifolds and Einstein manifolds
  of dimension four}, Compositio Math. 49 (1983), Proposition 5 and
  Theorem 2:
  \url{https://www.numdam.org/article/CM_1983__49_3_405_0.pdf}.
\item
  M. T. Anderson, \emph{Orbifold compactness for spaces of Riemannian
  metrics and applications}, including the recalled Einstein compactness
  case: \url{https://arxiv.org/abs/math/0312111}.
\item
  E. P. Wright, \emph{Quotients of gravitational instantons}, Theorem A:
  \url{https://arxiv.org/abs/1102.2442}.
\item
  A. Preissmann, \emph{Quelques propriétés globales des espaces de
  Riemann}, Comment. Math. Helv. 15, 175--216, original scan:
  \url{https://www.e-periodica.ch/cntmng?pid=com-001\%3A1942\%3A15\%3A\%3A15}.
\item
  X. Dai, X. Wang, G. Wei, \emph{On the Variational Stability of
  Kähler-Einstein Metrics}, Theorem 3.6 in the inspected author
  manuscript: \url{https://web.math.ucsb.edu/~dai/paper/dww3.pdf}.
\item
  P.-A. Nagy, second-order and third-order negative KE deformation
  results: \url{https://arxiv.org/abs/2603.09028} and
  \url{https://arxiv.org/abs/2606.04501}.
\end{itemize}

The application uses the complete hypotheses of the inherited shear
topology and actual monopole results. The new content of this note is
their composition with the spectral structure estimate to obtain the
global escape alternative, including the explicit zero-budget-ratio
application. It introduces neither a new compactness theorem nor a
general Einstein nonexistence result.

\end{document}
